\section*{الفصل \ref{ch:b3:x-ray-diffraction} --- علم البلورات وحيود الأشعة السينية}

\begin{solution}{exo:b3:x-ray-diffraction:1}
$d = a/\sqrt N$: $d_{111} = \qty{325.66}{pm}$، $d_{200} = \qty{282.03}{pm}$، $d_{220} =
\qty{199.42}{pm}$.
\end{solution}

\begin{solution}{exo:b3:x-ray-diffraction:2}
مكعب ممركز الأوجه: الانعكاسان الأولان هما (111) و(200). $\sin\theta = \lambda\sqrt N/2a$: $2\theta =
\ang{43.32}$ و\ang{50.45}.
\end{solution}

\begin{solution}{exo:b3:x-ray-diffraction:3}
P: السبعة كلها. I ($h + k + l$ زوجي): 110، 200، 211، 220. F (غير مختلطة): 111، 200، 220.
\end{solution}

\begin{solution}{exo:b3:x-ray-diffraction:4}
$\mathbf a^* = \mathbf e_x/a$، $\mathbf b^* = \mathbf e_y/b$: شبكة مستطيلة ضلعاها $1/a$
و$1/b$، والضلع الطويل في الشبكة المباشرة يصير القصير.
\end{solution}

\begin{solution}{exo:b3:x-ray-diffraction:5}
نسب $\sin^2\theta$ هي $1 : 2 : 3 : 4 : 5$، أي $N = 2, 4, 6, 8, 10$: ممركزة الجسم (110،
200، 211، 220، 310). ومن الخط الأول، $d_{110} = \lambda/2\sin(\ang{20.135}) =
\qty{223.77}{pm}$ و$a = d\sqrt2 = \qty{316.5}{pm}$: إنه التنغستن.
\end{solution}

\begin{solution}{exo:b3:x-ray-diffraction:6}
$Z = \rho N_Aa^3/M = 1.99 \times 6.022\times10^{23} \times (6.2929\times10^{-8})^3/74.6 = 4.0$.
\end{solution}

\begin{solution}{exo:b3:x-ray-diffraction:7}
$\beta = \qty{0.40}{\degree} = \qty{6.98e-3}{rad}$، $\cos\theta = \cos\ang{15.85} = 0.962$:
$L = 0.9 \times 0.15406/(6.98\times10^{-3} \times 0.962) = \qty{21}{nm}$.
\end{solution}

\begin{solution}{exo:b3:x-ray-diffraction:8}
$F_{hkl} = f_{\ce{Cs+}} + f_{\ce{Cl-}}(-1)^{h+k+l}$: $F_{100} = 54 - 18 = 36$، $F_{110} = 72$:
فالشدتان بنسبة 1 : 4. يشغل المركزَ أيونٌ مختلف عن أيونات
الرؤوس: فالشبكة أولية (فالشبكة الممركزة الجسم تقتضي محتوى متطابقًا
في النقطتين)، والانعكاس 100 حاضر.
\end{solution}

\begin{solution}{exo:b3:x-ray-diffraction:9}
لكل ذرة نسخة مزاحة بمقدار $(\frac12,\frac12,0)$، وهذا يضرب $F$ في $1 +
\eu^{\iu\pi(h+k)} = 1 + (-1)^{h+k}$، المعدوم من أجل $h + k$ فردي.
\end{solution}

\begin{solution}{exo:b3:x-ray-diffraction:10}
لأشباه البلورات ترتيب بعيد المدى دون دورية: فلا
شبكة لها، ولا ينطبق عليها القيد، الذي يخص الشبكات.
وأدت بقع حيودها الحادة إلى تعريف البلورة بأنها كل صلب له نمط
حيود منفصل في جوهره، دوريًا كان أم لا.
\end{solution}

\begin{solution}{exo:b3:x-ray-diffraction:11}
ينقل الانزلاق $c$ النقطة $(x, y, z)$ إلى $(x, \frac12 - y, \frac12 + z)$. ومن أجل $k = 0$ يكون عاملا
الطور $\eu^{2\pi\iu(hx + lz)}$ و$\eu^{2\pi\iu(hx + lz + l/2)} = (-1)^l
\eu^{2\pi\iu(hx+lz)}$: ومجموعهما ينعدم من أجل $l$ فردي.
\end{solution}

\begin{solution}{exo:b3:x-ray-diffraction:12}
الانقلاب أو المرآة أو الانزلاق يحوّل الجزيء إلى خياله المرآوي، فالبلورة
التي تحتوي أحدها تحتوي المتماكبين المرآويين كليهما. والمتماكب المرآوي النقي لا
يتبلور إلا في الزمر الفضائية الخمس والستين المبنية من الدورانات والمحاور اللولبية
والانسحابات (زمر سونكه)، مثل \babelsublr{\textit{P}$2_12_12_1$}. وبحسب قانون فريدل يكون
نمط أحد المتماكبين المرآويين مساويًا لنمط الآخر؛ ويُستخرج الترتيب الفراغي
المطلق من الانحرافات الصغيرة التي يسببها الانتثار الشاذ
للذرات الأثقل.
\end{solution}

\begin{solution}{pb:b3:x-ray-diffraction:1}
\textbf{1.} $\theta = \ang{14.171}$، $d = \lambda/2\sin\theta = \qty{314.64}{pm}$.
\textbf{2.} 314.64، 222.49، 181.66، 157.32، 140.71، \qty{128.45}{pm}.
\textbf{3.} 1.000، 2.000، 3.000، 4.000، 5.000، 6.000.
\textbf{4.} 1، 2، 3، 4، 5، 6.
\textbf{5.} نعم، بوصفها (100) و(110) و(111) و(200) و(210) و(211) لخلية مكعبة أولية
حيث $a' = \qty{314.6}{pm}$.
\textbf{6.} $N = 7$ (و15 و23\dots)، وهو ليس مجموع ثلاثة مربعات؛ 
وأول فرق سيظهر عند الخط الثامن. فستة خطوط لا تكفي للتمييز بين خلية P
حرفها $a'$ وخلية F حرفها $2a'$.
\textbf{7.} $N = 4, 8, 12, 16, 20, 24$: (200)، (220)، (222)، (400)، (420)، (422).
\textbf{8.} $a = 2d_{200} = \qty{629.3}{pm}$.
\textbf{9.} \ce{KCl} (\qty{629.29}{pm}).
\textbf{10.} $F = 4[f(\ce{K+}) - f(\ce{Cl-})]$ من أجل القرائن الفردية كلها، ولكل من الأيونين ثمانية عشر
إلكترونًا: فتتلاشى السعتان.
\textbf{11.} نعم في الحالتين: \ce{Na+} (10) و\ce{Cl-} (18)، و\ce{K+} (18) و\ce{Br-} (36)
تنثر انتثارًا مختلفًا؛ فالانعكاس (111) للمركب \ce{NaCl} يقع عند \ang{27.36}، وللمركب \ce{KBr} عند \ang{23.38}،
وكلاهما داخل المجال المسجَّل.
\textbf{12.} ترى الأشعة السينية \ce{K+} و\ce{Cl-} شبه متطابقين؛ والأيونات المتطابقة في
ترتيب الملح الصخري تشكّل مصفوفة مكعبة بسيطة حرفها $a/2$، وهي خلية ظاهرية.
\textbf{13.} أربع وحدات \ce{KCl}.
\textbf{14.} $\rho = 4 \times 74.6/(6.022\times10^{23} \times (6.293\times10^{-8})^3) =
\qty{1.99}{g/cm^3}$.
\textbf{15.} ستة وستة (ثمانيات أوجه من الأيون الآخر).
\textbf{16.} $a/2 = \qty{314.6}{pm}$.
\textbf{17.} (440)، $N = 32$، عند \ang{87.65} (والانعكاس الفردي كله (333)/(511) غائب).
\textbf{18.} تتغير الشدات مع الاتجاه المفضَّل للبليرات، 
والامتصاص، وتناقص عوامل الانتثار والحركة الحرارية؛ أما المواضع
فلا تتعلق إلا بالشبكة.
\textbf{19.} من $d = \lambda/2\sin\theta$، $\Delta d/d = -\cot\theta\,\Delta\theta$: فالخطأ الزاوي نفسه
يعطي خطأً نسبيًا أصغر عند $\theta$ الكبيرة.
\textbf{20.} $d_{420} = \lambda/2\sin\ang{33.190} = \qty{140.71}{pm}$، $a = d\sqrt{20} =
\qty{629.3}{pm}$.
\textbf{21.} $L = 0.9 \times 0.15406/(4.36\times10^{-3} \times \cos\ang{33.19}) = \qty{38}{nm}$.
\textbf{22.} لا: $\beta = \ang{0.0095}$، وهو أدنى بكثير من العرض الآلي النموذجي.
\textbf{23.} \textbf{$a = \qty{629.3}{pm}$: المسحوق هو كلوريد البوتاسيوم}، وقد أسكت انعكاساتِه
الفرديةَ أيونان لهما العدد نفسه من الإلكترونات.
\end{solution}
